
\documentclass{article}
\usepackage{colortbl}
\usepackage{makecell}
\usepackage{multirow}
\usepackage{supertabular}

\begin{document}

\newcounter{utterance}

\centering \large Interaction Transcript for game `cladder', experiment `full\_v1.5\_default', episode 188 with Qwen3.5{-}27B.
\vspace{24pt}

{ \footnotesize  \setcounter{utterance}{1}
\setlength{\tabcolsep}{0pt}
\begin{supertabular}{c@{$\;$}|p{.15\linewidth}@{}p{.15\linewidth}p{.15\linewidth}p{.15\linewidth}p{.15\linewidth}p{.15\linewidth}}
    \# & \multicolumn{2}{c}{Player} && \multicolumn{2}{c}{Game Master} \\
    \hline

    \theutterance \stepcounter{utterance}  
    & & & \multicolumn{4}{p{0.6\linewidth}}{
        \cellcolor[rgb]{0.9,0.9,0.9}{
            \makecell[{{p{\linewidth}}}]{
                \texttt{\tiny{[P1$\langle$GM]}}
                \texttt{You are an expert in causal inference. The following question is not a typical commonsense query, but rather a meticulously designed question created by a professor specializing in causal inference, intended to assess the students' mastery of the course content.} \\
\texttt{Imagine a self{-}contained, hypothetical world with only the following conditions, and without any unmentioned factors or causal relationships: The captain has a direct effect on the private and the corporal. The corporal has a direct effect on prisoner. The private has a direct effect on prisoner. Method 1: We look directly at how the captain correlates with prisoner in general. Method 2: We look at this correlation case by case according to the private. To understand how the captain affects prisoner, is it more correct to use the Method 1 than Method 2?} \\
\texttt{Start your answer with 'Yes' or 'No', followed by additional reasoning or evidence to support your explanation.} \\
            }
        }
    }
     \\ \\

    \theutterance \stepcounter{utterance}  
    & \multicolumn{4}{p{0.6\linewidth}}{
        \cellcolor[rgb]{0.561,0.737,0.561}{
            \makecell[{{p{\linewidth}}}]{
                \texttt{\tiny{[P1$\rangle$GM]}}
                \texttt{No} \\
\\ 
\texttt{In this hypothetical world, the causal structure is defined as: Captain \$\textbackslash{}to\$ Private, Captain \$\textbackslash{}to\$ Corporal, Private \$\textbackslash{}to\$ Prisoner, and Corporal \$\textbackslash{}to\$ Prisoner. This creates a Directed Acyclic Graph (DAG) where the Private acts as a **mediator** on the causal path from the Captain to the Prisoner (Captain \$\textbackslash{}to\$ Private \$\textbackslash{}to\$ Prisoner).} \\
\\ 
\texttt{Method 1 (looking at the general correlation) estimates the **Total Effect** of the Captain on the Prisoner. This includes both the direct effect (via the Corporal) and the indirect effect (via the Private). Since the goal is to understand "how the captain affects prisoner" in a general sense, capturing the total impact is the standard objective in causal inference unless a specific direct effect is requested.} \\
\\ 
\texttt{Method 2 (stratifying or conditioning on the Private) blocks the causal path that goes through the Private. By conditioning on a mediator, you effectively "cut" the indirect pathway, isolating only the direct effect of the Captain on the Prisoner (via the Corporal). This results in an estimate that excludes a significant portion of the Captain's actual influence on the Prisoner. In causal inference, conditioning on a mediator to estimate the total effect is a form of **over{-}control bias** (or simply estimating the wrong quantity), as it removes the very mechanism through which part of the cause operates.} \\
\\ 
\texttt{Therefore, Method 1 is more correct for understanding the overall effect of the Captain on the Prisoner, whereas Method 2 would yield a biased estimate of the total effect by artificially removing the indirect pathway.} \\
            }
        }
    }
    & & \\ \\

    \theutterance \stepcounter{utterance}  
    & & & \multicolumn{2}{p{0.3\linewidth}}{
        \cellcolor[rgb]{0.9,0.9,0.9}{
            \makecell[{{p{\linewidth}}}]{
                \texttt{\tiny{[GM$|$GM]}}
                \texttt{No} \\
            }
        }
    }
    & & \\ \\

    \theutterance \stepcounter{utterance}  
    & & & \multicolumn{2}{p{0.3\linewidth}}{
        \cellcolor[rgb]{0.9,0.9,0.9}{
            \makecell[{{p{\linewidth}}}]{
                \texttt{\tiny{[GM$|$GM]}}
                \texttt{yes} \\
            }
        }
    }
    & & \\ \\

    \theutterance \stepcounter{utterance}  
    & & & \multicolumn{2}{p{0.3\linewidth}}{
        \cellcolor[rgb]{0.9,0.9,0.9}{
            \makecell[{{p{\linewidth}}}]{
                \texttt{\tiny{[GM$|$GM]}}
                \texttt{game\_result = LOSE} \\
            }
        }
    }
    & & \\ \\

\end{supertabular}
}

\end{document}
